\documentclass[10pt]{article}

\usepackage[utf8]{inputenc}

\usepackage[T1]{fontenc}

\usepackage{amsmath}

\usepackage{amsfonts}

\usepackage{amssymb}

\usepackage[version=4]{mhchem}

\usepackage{stmaryrd}

\usepackage{bbold}

\usepackage{mathrsfs}



\begin{document}

ство коміактно (хотя и не обязательно отделимо). Теория У. Л. групп класса $[M A P]$ связана с теорией почти периодич, функций на локально компактных группах.



$\mathrm{X}$ a p a к т е р о м У. п. локально компактной групшы $G$ наз. такой точный нормальный полуконечный след $t$ на множестве положительных әлементов алгебры Неймана $\mu_{\pi}$, порожденной семейством $\pi(G)$, что множество элементов $x$ из групповой $C^{*}$-алгебры $C^{*}(G)$ грушиы $G$ (обертывающей $C^{*}$-алгебры групповой алгебры $\left.L_{1}(G)\right)$ таких, что $t\left(\pi\left(x^{*}, x\right)\right)$ конечен, переходит в множество, порождающее алгебру Неймана $\boldsymbol{\pi}(G)^{\prime \prime}$. Если $\pi$ - факторпредставление (соответственно неприводимое У. п.), то характер $t$ определяет У. ІІ. $\pi$ однозначно с точностью до квазиәквивалентности (соответственно эквивалентности). Если $\pi_{1}, \pi_{2}$ - неприводимые У. п. группы $G$ с характерами $t_{1}, t_{2}$ соответственно, то произведение этих следов определяет след на алгебре Неймана, пгорожденной У . п. $\pi_{1} \otimes \pi_{2}$. Если этот след является характером представления $\pi_{1} \otimes \pi_{2}$, то он (для сепарабельной групшы или сепарабельных пространств представлений) определяет разложение У. п. $\pi_{1} \bigotimes \pi_{2}$ в прямої интеграл факторпредставлениӥ со следом по однозначно определенной мере (мере Планшереля для $\pi_{1}\left(\pi_{2}\right.$ ) на квазиспектре $\hat{G}$. Нахождение мер Планшереля тензорных произведений У. п. является одной из общих задач теории У. П.; в ряде случаев (в частности, для групп $S L(2, \mathbb{C}), S L(2, \mathbb{R})$, нек-рых У. п. др. полупростых групп Ли и нек-рых разрешимых групп Ли) эта задача рептена (с помощыо изучения спектрального разложения оператора Лапласа, методом орбит или методом орисфер).



Иногда х а р а к т е р о м У. п. $\pi$ в гильбертовом пространстве $H$ наз. линейный функционал $\chi \pi$ на инвариантной относительно сдвигов Іодалгебре $D_{\pi}(G)$ в алгебре $M(G)$, определенный равенством $\chi \pi(a)=\operatorname{Tr} \tilde{\pi}(a)$, $a \in D_{\pi}(G)$, где $\tilde{\pi}-$ представление алгебры $D_{\pi}(G)$, определенное $\pi$ (гредполагается, что $\pi$ однозначно определяется шредставлением $\tilde{\pi}$, онераторы представления $\tilde{\pi}$ ядерны и отображение $\tilde{\pi}$ алгебры $D_{\pi}(G)$ в пространство ядерных операторов непрерывно). Характеры неприводимых У. п. полупростых и нильпотентных групп Ли определяются обобщенными функциями, к-рые в случае полупростых групп измеримы и локально интегрируемы. Характеры неприводимых У. п. разрешимых груип Ли типа I определены, вообще говоря, лишь на подалгебрах алгебры $C_{0}^{\infty}(G)$ финитных бесконечно дифференцируемых функций на $G$. Вычисление характеров во многом основано на формуле для характера индуцированного представления.



Ћомпактная подгрупта $K$ группы $G$ наз. м а с с и вн о й, если ограничение любого неприводимого У. г. $\pi$ 'группы $G$ на подгруппу $K$ содержит любое неприводимое У. п. $\sigma$ группы $K$ с конечной кратностью. Пусть $P_{\sigma}^{\pi}-$ проектор в пространстве $H_{\pi}$ представления $\pi$ на подпространство, в к-ром действует представление групшы $K$, кратное $\sigma$; функции вида



\[

\varphi_{\boldsymbol{\sigma}}^{\pi}(g)=\operatorname{Tr}\left(P_{\boldsymbol{\sigma}}^{\pi} \pi(g)\right), \quad g \in G,

\]



наз. $K$-с ф е ритеск ими фу укция ми представления $\pi$ (ср. Представляющая функция). Группа $G$ с массивной компактной подгруппой принадлежит типу I; каждое неприводимое У. п. группы $G$ имеет характер и определяется однозначно с точностью до эквивалентности любой ненулевой сферич. функцией; дуальное пространство $\hat{G}$ к группе $G$ можно представить в виде счетного объединения (пересекающихся) хаусдорфовых локально компактных пространств (образованных такими $\pi \in \widehat{G}$, что $\varphi_{\sigma}^{\pi}=0$ для данного $\sigma \in \hat{K}_{\pi}$, а размерность соответствующего проектора $P_{\sigma}^{\pi}$ имеет данное значение). К числу групп с массивной компактной подгруюиой относятся линейные нолупростые группы Ли и [Z]-групшы.



Пусть $\pi-У$. П. группы $G$ в гильбертовом пространстве $H, M_{\pi}$ - алгебра Неймана, порожденная семейством $\pi(G)$. Представление $\pi$ наз. д о і у с к а ю щ и м с л е д, если существует след на $M_{\pi}^{+}$, являющийся характером У. п. л. С л е д о м н а Іолуконечныйі полунепрерывный снизу след на $C^{*}(G)^{+}$; след на $G$ наз. х а р а к т е р о м группы $G$, если соответствующее У. и. групиы $G$ есть факторпредставление. Существует каноническое взаимно однозначное соответствие между множеством характеров группы $G$, определенных с точность до положительного множителя, и множеством классов квазиэквивалентности факторпредставлений группы $G$, допускающих след; при этом факторпредставлениям конечного типа соответствуют непрерывные центральные положительно огределенные функции на $G$.



Реәуланое представление локально компактной группы $G$ в гильбертовом пространстве $L_{2}(G)$ есть точное непрерывное У. п.; $C^{*}$-алгебра, порожденная образом соответствующего гредставления алгебры $L_{1}(G)$, наз. п р и в е д е н о о ї $C^{*}$-а л г е б р о й груіпы $G$ и обозначается $C_{r}^{*}(G)$; пусть $N$-ядро канонического эпиморфизма $C^{*}(G)$ на $C_{r}^{*}(G)$, определяемого регулярным представлением. Груіна $G$ а м е н а б е л ь н а, т. е. на $L_{\infty}(G)$ существует инвариантное среднее, тогда и только тогда, когда $N=\{0\}$ (ограниченное представление аменабельной группы в гильбертовом пространстве әквивалентно У. П.). Семейство таких У. п. $\pi \in \hat{G}$, что ядро соответствующего представления $C^{*}$-алгебры $C^{*}(G)$ содержит $N$, наз. о с н о в н о й с е р и е й; остальные У. П. $\pi \in \widehat{G}$ образуют д о по ли и т е льн у ю с е р и ю.



Пусть $G$ - унимодулярная сепарабельная локально компактная группа типа I, $W^{*}(G)$ - алгебра Неймана $\mu_{\lambda}$, где $\lambda$ - регулярное У. п. группы $G$. Существует едцинственная положительная мера $\bar{\mu}$ на спектре $\hat{G}$ группы $G$, удовлетворяющая условию



\[

\int_{G}|f(g)|^{2} d \mu(g)=\int_{\hat{G}} \operatorname{Tr}\left(\pi(f) \pi(f)^{*}\right) d \hat{\mu}(\pi)

\]



для всех $f \in L_{1}(G) \cap L_{2}(G)$. Мера $\hat{\mu}$ наз. м е р о ї П л а нше р е л. Кроме того, существует изоморфизм пространства $L_{2}(G)$ на прямой интеграл операторов Гильберта - Шмидта в гространствах представлений $\pi \in \widehat{G}$ по мере $\hat{\mu}$, переводящий левое регулярное У. П. $\lambda$ в прямой интеграл У. І., кратных представлениям $\pi$, а след на $W^{*}(G)$, определяемый следом $\mathscr{E}_{e}$ на $G\left(\mathscr{E}_{e}(f)=\right.$ $=f(e)$ для $f \in K(G)$, в прямой интеграл следов $T \otimes 1 \rightarrow$ $\rightarrow \operatorname{Tr} T$ на $\left(\mu_{\pi} \otimes \mathbb{C}_{1}\right)^{+}$. След $\mathscr{E}_{e}$ на $C^{*}(G)+$ совпадает со следом $f \rightarrow \int_{\hat{G}} \operatorname{Tr} \pi(f) d \hat{\mu}(\pi), f \in C^{*}(G)$. Формула (:) наз. ф о р м у л о й П л а н пе е е ля; она допускает обобщение на несепарабельные унимодулярные локально компактные группы типа I, а также на неунимодулярные сепарабельные локально компактные группы и сепарабельные группы не типа І. Одной из задач теории У. п. является явное построение меры Планшереля для данной локально компактной группы. Эта задача решена лишь частично (напр., для полупростых вещественных групп Ли, разрешимых групп Ли типа I, а также некоторых групп движений, некоторых групп ШІевалле и нек-рых групп с условиями компактности). С разложением регулярного У. п. и формулой Планшереля связана теория квадратично интегрируемых представлений, дискретных серий и интегрируемых представлений.



Полное описание неприводимых У. П. локально компактных групі даже в случае групп Ли не известно.





\end{document}